\chapter{Acceleration, Strain, and Curvature}
\label{chap:finite-extrapolation-alpha}

\section{After the Loop}

The previous chapter brought the electron to the Cooper-pair boundary. Charge
was no longer an open motion. It was a loop count. Mass was no longer a hidden
substance. It was second-difference strain. Phase was not decoration. It was
the cost of keeping orientation coherent. Sameness was not free. It was the
bounded quotient by which the paired transport could stand as one unit for a
specified reading.

Chapter 09 begins from that boundary. The grammar now asks how a finite reader
can extract the fine-structure reading, \(\alpha\), without pretending that a
completed constant has been handed down from outside the trial. The answer
cannot be a direct declaration of a number. A direct declaration would skip the
ledger that made charge, mass, phase, and sameness readable. The chapter must
show how the finite record carries enough structure to extrapolate \(\alpha\)
while still carrying the residue left by every bounded trial.

The word extrapolate is important. A measurement does not possess the limiting
value merely because the notation can name it. A finite trial samples an
aperture, records a bounded result, and leaves a residue. Another trial may
refine the reading, but it must inherit the state left by the previous one.
If the new trial restarts the record, it is not a refinement of the old trial.
It is a new question with a similar name. The grammar therefore permits
extension and forbids amnesia.

The chapter follows four obligations. First, a repeated trial must carry seed
and residue without restarting. The seed is the selector by which the next
finite sample is placed in the trial. The residue is the unfinished remainder
that the previous sample could not absorb. Together they make repetition a
continuation rather than a stack of unrelated attempts.

Second, the aperture must be read through three rungs. The same bounded
opening can be faced as charge, mass, or value. Charge reads loop count and
orientation. Mass reads second-difference strain. Value reads the normalized
result that can be compared across trials. A fourth face is not a fourth kind.
It returns to charge by the same corridor, and the grammar must quotient that
return rather than multiply the rungs.

Third, the finite readings must be extrapolated by a bounded Richardson-style
move. A depth is sampled. The depth is doubled. The leading sampling error is
read from the disagreement between the two finite depths. Only when that error
is forced to change side inside a bounded bracket may the grammar improve the
estimate. The extrapolated value is not omniscience. It is the reading trapped
by two finite depths and an explicitly carried remainder.

Fourth, the path relation must separate signed displacement from cost. A
direct path and a composed path may carry the same signed displacement. The
grammar may identify that displacement when the residues compose. But the two
paths need not cost the same. The composed path spends an intermediate
transition. The direct path bypasses it. Alpha is not safely recovered until
the ledger has accounted for both the signed displacement and the cost by
which the displacement was obtained.

Examples can illuminate these moves. Repeated observations cluster only when
the alphabet and procedure stay fixed. A laser tracking record grows by
emission and return events instead of restarting at each pulse. A coastline
length changes under refinement unless a residual law says what refinement is
allowed to preserve. A slow boundary can dominate the cost of reconciliation
even when the displacement is short. Such examples are useful lanterns, but
they are not the spine of the chapter.

The spine is grammatical. The grammar permits the finite trial, carries the
seed, restores the residue, reads the aperture through three rungs, bounds the
depth-doubling extrapolation, and keeps path cost distinct from signed
displacement. Under those obligations, \(\alpha\) can enter as a finite
extrapolated reading. It is neither a primitive constant nor a loose empirical
average. It is the dimensionless ratio forced when the carried loop, the
strain face, the value rung, and the bounded path relation are made answerable
to one ledger.

\section{The Accelerated Frame}

A repeated trial is not a row of fresh beginnings. If every pass begins from
zero, the passes may resemble one another, but they do not refine one another.
They are separate questions. The grammar of alpha cannot use that kind of
repetition, because \(\alpha\) is not read from a single isolated sample. It
is read from a bounded process in which each sample inherits what the previous
sample did not settle.

The inherited state has two essential pieces: seed and residue. The seed is
the carried selector by which the next sample is placed. It tells the grammar
which admissible face of the aperture is current. It is not an unaccountable
whim and not an appeal to randomness as explanation. It is the part of the
trial state that makes the next selection repeatable as a successor of the
old one. Without a carried seed, the next sample may be possible, but it is
not yet the next sample of this trial.

The residue is the unfinished remainder. It records the fraction of a turn,
the leftover strain, or the unabsorbed reading that has not yet crossed the
trial's counting boundary. A residue that disappears between trials is not a
residue; it is a discarded error. The grammar forbids that discard. If a
finite pass leaves an unread remainder, the next pass must begin with that
remainder in the ledger.

This is why the trial advances by carrying rather than restarting. A pass
selects an aperture face, samples the finite neighborhood allowed by that
face, projects the mean remainder onto the current lattice of the trial, and
adds it to the residue already present. If the sum crosses a full turn, the
grammar counts the completed turn. If it does not, the grammar carries the
remainder. In either case, the new state is not blank. It is the old state
extended.

The wrap is not a loss of information. When a residue crosses one full turn,
the completed turn is counted and the remainder is restored as the starting
residue for the next pass. The count and the residue therefore separate two
faces of the same event. The count says what has closed. The residue says
what remains open. A trial that keeps only the count loses the local memory
needed for the next aperture. A trial that keeps only the residue loses the
closed loop that makes charge billable.

The seed is carried through the same extension. Once a face has been sampled,
the selector advances. The next face is not chosen from outside the record.
It is obtained by the successor state of the trial. The grammar may therefore
read the next sample as belonging to the same repeated process. A stable
procedure does not mean that the same face must be sampled every time. It
means that the rule by which faces are sampled remains part of the record.

This distinction protects the trial from a common statistical shortcut.
Repeated measurements can be averaged only when the alphabet of measurement
and the update rule have remained stable. If the alphabet changes between
passes, the mean is a mixture of incompatible ledgers. If the update rule
changes, the spread no longer reports the residue of one procedure. The
grammar of alpha requires a stronger condition: not merely many samples, but
many samples whose seed and residue form one bounded history.

The finite bound is equally important. A residue carried without a horizon
would become an excuse to defer every debt. The grammar must know how much
of the trial has been spent, how many passes have been permitted, and what
remainder is present when the bound is reached. A bounded trial may stop with
residue still visible. That is not failure. It is the honest statement of
what the trial has read and what it has not read.

The laser-tracking example gives a clean image. A reader who emits a pulse
and records its return has not erased the previous pulse. Each emission and
return pair constrains the path already being tracked. A long record is not
many beginnings; it is one growing chain of commitments. If the tracked body
accelerates, later returns force new constraints into the same ledger. The
chain grows because the earlier constraints remain in force.

The same form governs invisible motion between anchors. If two readers refine
an interval and no new event appears, the grammar does not say that the
interval was empty. It says that the carried residue has been driven below
the current distinguishability threshold for that trial. If a later shared
anchor exposes a disagreement in value, slope, or bending moment, the
disagreement becomes residue that must be carried forward. Repeatability is
the rule that makes such residue communicable.

For alpha, the carried residue has a special consequence. A single aperture
sample cannot decide the fine-structure reading, because the reading belongs
to the relation between loop count, strain, and value. Each pass samples one
face under one finite selection. The residue of that pass supplies the next
pass with its starting debt. The seed supplies the next pass with its
admissible selector. The trial becomes an orbit of bounded obligations.

The grammar therefore reads finite repetition in three layers. It permits a
next pass only as an extension of the prior state. It carries seed and residue
as the state that makes the next pass continuous with the last. It bounds the
sequence so the carried residue remains a readable debt rather than a promise
of future absolution. Only under those rules can repeated trials contribute
to a finite extrapolation of \(\alpha\).

The next question is what the aperture itself can read. The carried trial
does not sample a formless opening. It samples a bracket with three faces.
Those faces are the rungs on which the alpha reading will stand: charge,
mass, and value.

\section{Second Difference as Strain}

The aperture is a bracket before it is a value. It has a lower side, an upper
side, and a width the grammar can read. The bracket says that the current
trial has trapped a live region: too low has not yet reached the boundary,
too high has passed it, and the aperture between them is the finite opening
inside which the next sample must be placed. The grammar permits the bracket
to be refined. It forbids the bracket from being treated as though every
interior point had already been measured.

The same aperture can be read through three rungs:
\[
  \hbox{charge},\qquad \hbox{mass},\qquad \hbox{value}.
\]
These rungs are not three independent constants. They are three faces of the
same bounded opening. The trial may select one face at a pass, but the bracket
itself remains the shared aperture. What changes is the question being asked
of the bracket: how many signed turns have closed, how much strain the
closure carries, or what normalized result can be compared with the next
trial.

The charge rung reads loop count and orientation. It asks whether the carried
residue has crossed a full turn and whether that completed turn faces the
reader with the admitted sign. The charge rung inherits Chapter 08 directly.
Charge is not a substance moving through the aperture. It is the counted
closure of a loop whose sign has become accountable. When the aperture is
read as charge, the grammar reads the bracket as a place where a signed
turn may be counted or withheld.

The mass rung reads strain. It asks how the closed loop responds when the
trial has passed null and threshold into a second difference. The aperture
does not merely say that a turn has occurred. It asks whether the response
bends. A mass reading is therefore the strain face of the same bracket that
charge read as count. If the response carries no second-difference debt at
this scale, the mass rung has no strain to report. If it does, the strain
enters the ledger as cost.

The value rung reads the normalized result. It is the face by which the
trial can compare one bounded reading with another. Value does not mean that
the bracket has become a completed real number. It means that the bracket
has supplied a result stable enough to be carried across the next comparison.
The value rung is where the trial can ask whether the carried charge and
strain have yielded the same dimensionless reading after the next refinement.

Thus the aperture is one opening with three admissible readings. The grammar
does not let charge forget mass, because charge without strain cannot scale
the response. It does not let mass forget charge, because strain without
loop count is a cost with no closed traversal to bill. It does not let value
forget either one, because a normalized result detached from count and strain
would be only a number-name. The alpha reading requires all three faces to
belong to one finite aperture.

The order of the rungs matters because the grammar is cyclic:
\[
  \hbox{charge} \to \hbox{mass} \to \hbox{value} \to \hbox{charge}.
\]
The fourth reading is not a new rung. It returns to the first. This return
is the aperture's local version of the loop coming home. A finite aperture
can resolve charge, mass, and value; when a fourth face is demanded, the
grammar identifies it with charge under the corridor quotient. Without that
quotient, the trial would manufacture an unlimited ladder of names from a
three-rung bracket.

This ceiling is not poverty. It is the condition that makes the reading
auditable. A trial that could invent a new face whenever the old three became
inconvenient would never force a result. It would evade residue by changing
the vocabulary. The three-rung aperture forbids that evasion. Every pass must
ask a charge question, a mass question, or a value question. Every return
after value must face charge again.

The precision example clarifies the value rung. A spline completion can
assign a determinate value inside a bracket once anchors and a completion
rule have been fixed. That value is not an independent measurement. It is a
consequence of the bracket and rule. In the same way, the value rung of
\(\alpha\) supplies a determinate finite result only because charge and mass
have already constrained the aperture. Future trials may refine it, but they
must refine the same carried obligation.

The aperture also keeps the statistical reading honest. A mean or variance
computed over incompatible faces would mix ledgers. The charge face has one
moment structure, the mass face another, and the value face another. The
grammar permits comparison only when the face has been named and the bracket
has been preserved. The same lower and upper aperture can support all three
faces, but the selected face tells the trial which moment is currently being
read.

This is where \(\alpha\) becomes more than a charge receipt. Alpha is not the
elementary charge itself. Nor is it the mass strain alone. It is the finite
dimensionless reading that appears when the charge rung, the mass rung, and
the value rung are brought into one bounded comparison. The charge supplies
the signed loop. The mass supplies the strain scale. The value supplies the
carried ratio. The aperture forces them to answer together.

The three-rung rule also explains why the chapter cannot jump directly to a
final numerical constant. A number without the aperture would not know which
face it had read. A face without the bracket would not know its bound. A
bracket without the cyclic return would let value drift away from charge and
mass. The grammar therefore reads \(\alpha\) only after the aperture has
carried all three rungs and restored the cycle.

Once the aperture has been bounded and runged, refinement can begin. But
refinement is dangerous. A smaller aperture is not automatically a better
reading. The grammar must say how a coarse depth and a doubled depth can be
compared, how their sampling errors are trapped, and why the extrapolated
reading remains bounded. That is the Richardson-style obligation.

\section{The Ricci Reading}

Refinement is not automatically improvement. A finer trial may expose a
residue that the coarser trial could not see. It may also spend more cost
without narrowing the reading that matters. The grammar therefore cannot say
``run deeper'' and call the result knowledge. It must specify which deeper
trial refines the old one, how the old residue is carried, and what error is
being cancelled or bounded.

The Richardson-style move begins with two finite depths. At depth \(d\), the
trial reads an aperture value \(a_d\). At doubled depth \(2d\), it reads
another value \(a_{2d}\). Both readings are finite. Both are bounded. The
doubled reading is not a completed limit. It is a second trial that must
preserve the seed, residue, and aperture language of the first trial while
sampling the bracket more tightly.

The grammar licenses extrapolation only when the leading sampling error can
be read with a sign. A coarse depth may read the value from one side of the
aperture. The doubled depth may force the leading error to the other side.
When that happens, the disagreement between \(a_d\) and \(a_{2d}\) is not
mere noise. It is a signed witness. The error has faced the reader, crossed
the bracket, and become eligible for cancellation.

One schematic form is enough:
\[
  a_d = \alpha + e_d,
  \qquad
  a_{2d} = \alpha - \lambda e_d + r_d,
  \qquad
  0 \leq |r_d| \leq \epsilon_d .
\]
The notation records the grammar rather than founding it. The first reading
carries a leading signed error \(e_d\). The doubled reading carries the same
error with opposite face and controlled scale \(\lambda\), plus a remainder
\(r_d\). The bound \(\epsilon_d\) says how much unread residue remains after
the paired-depth comparison.

From those two finite readings the grammar may form the improved estimate
\[
  \widehat{\alpha}_{2d}
    = \frac{a_{2d}+\lambda a_d}{1+\lambda}.
\]
If the sign relation holds, the leading error cancels and the remaining
uncertainty is bounded by \(\epsilon_d/(1+\lambda)\). The chapter does not
need the formula as a ritual. It needs the obligation the formula records:
the extrapolated value is lawful only because two finite depths have trapped
the leading error between them and left a smaller explicit residue.

If the error does not change side, the extrapolation is not forced. Two
readings on the same side of the aperture may show drift, but they do not
authorize cancellation. The grammar may carry both readings, widen the
residue, or require another depth. It may not average them and pretend the
leading error has been removed. A signless disagreement is a warning, not a
Richardson certificate.

The doubling itself is a grammatical act. It is not the fantasy that every
intermediate point has been sampled. It is a finite rule: take the current
depth, refine by the agreed factor, and preserve the prior commitments unless
a new bounded residue explicitly revises them. If the doubled trial changes
the seed discipline, the aperture face, or the meaning of value, it is not a
Richardson partner. It is a different experiment.

The coastline example shows why this caution matters. A coast measured with
a shorter ruler may report a longer path because bays and inlets that were
previously bridged now enter the path. Refinement has changed the reading.
That change is not a mistake, but it also is not convergence by itself. To
extrapolate a boundary reading, the grammar needs a residual law that says
which roughness is being exposed, how the exposure scales with depth, and
why the remaining error is bounded.

The finite-precision example shows the opposite danger. Refinement can reach
a floor where new distinctions no longer enter as new events. Additional
operations then do not sharpen the value; they merely recycle existing
residue. A bounded extrapolation must know when it has reached such a floor.
The grammar therefore carries a machine-scale residue: the smallest
distinction the trial is authorized to spend at the current depth.

For \(\alpha\), the paired-depth reading has to preserve the three-rung
aperture. The charge rung must still read loop count and orientation. The
mass rung must still read strain. The value rung must still read the
normalized ratio. A doubled depth that sharpens only one rung while leaving
the others unaccounted for cannot force an alpha reading. It may be useful
data, but it is not yet the dimensionless coupling of the three rungs.

Depth doubling also carries the seed and residue from the previous section.
The second depth does not begin as if the first had never happened. Its seed
is the successor selector of the trial, and its residue includes the
unfinished remainder carried to the doubled sampling. This is why the
extrapolated value belongs to one process. It is not the agreement of two
strangers. It is the cancellation of a signed error inside one bounded
history.

The result is a finite extrapolation, not a completed revelation. The
grammar may say: these two depths, under this aperture and carried residue,
force the reading into this smaller bracket. It may not say: all future
depths have been consumed. The bound remains part of the result. A later
trial may refine the bracket again, but it must answer to the estimate and
residue already carried.

This is the Richardson obligation in the chapter. The grammar samples,
doubles, reads a sign-flipped error, cancels the leading term, bounds the
remainder, and carries the improved value forward as a finite estimate of
\(\alpha\). The next section asks what this estimate means along a path.
Two paths may deliver the same signed displacement through the rungs. They
may still spend different costs, and the alpha reading must keep that
difference visible.

\section{The Weyl Fence}

An extrapolated value is still a path reading. It has not floated free of the
route by which it was obtained. The grammar must therefore ask two questions
that ordinary shorthand often collapses. What signed displacement has the path
carried? And what cost did the path spend to carry it? The first question is
about orientation and endpoints. The second is about the ledger of transitions
used to get there.

The difference appears sharply in the three-rung corridor. Begin at the mass
rung and return to the charge rung. One route is direct: mass to the returned
charge face. Another route is composed: mass to value, then value to returned
charge. If the residues compose correctly, both routes carry the same signed
displacement through the corridor. The grammar may identify that displacement.
It may say that the endpoint relation has returned the same oriented result.

In schematic form:
\[
  \Delta_{\pm}(\gamma_{12}\circ\gamma_{23})
    = \Delta_{\pm}(\gamma_{13}) .
\]
The symbol \(\Delta_{\pm}\) names the signed displacement. The composed path
\(\gamma_{12}\circ\gamma_{23}\) passes through an intermediate rung. The
direct path \(\gamma_{13}\) does not. The equality says only that the signed
endpoint displacement is the same. It does not say that the two paths have
spent the same cost.

Cost is a different reading:
\[
  C(\gamma_{12}\circ\gamma_{23})
    \quad\hbox{and}\quad
  C(\gamma_{13}) .
\]
The composed path spends two transitions. It carries the intermediate value
reading, restores the residue at that middle face, and then carries the next
transition. The direct path bypasses the middle face. Even when the signed
displacement identifies, the costs may differ. The grammar forbids the cost
ledger from being quotiented merely because the displacement has been
quotiented.

This is not a minor bookkeeping point. It is what keeps alpha finite. If the
grammar erased cost whenever endpoint displacement agreed, then every cheaper
and more expensive route with the same sign would become the same trial. The
bounded Richardson estimate would lose its bound, because the route by which
the bound was obtained would no longer be readable. A finite extrapolation
must know both what displacement was carried and what work was spent to carry
it.

The composed path carries an intermediate residue. That residue may settle
inside the value rung before the returned charge face is read. The direct
path does not expose that middle residue. It may arrive at the same signed
endpoint because the corridor relation composes. But the absence of an
intermediate exposure is itself a difference in cost. The grammar therefore
identifies the displacement while preserving the route.

This is the path version of the earlier charge/mass distinction. Charge
counts closure. Mass reads strain. A single closed loop can have both faces,
but the count is not the strain. In the path relation, signed displacement
counts the oriented change through the corridor. Cost reads the transitions,
settling, and residue restoration by which that change was obtained. They
belong to the same path family, but they are not the same reading.

The tail-latency example shows the point at a boundary. A message may have a
short displacement in presentation and still cost more because it must be
reconciled across many connected interfaces. The slowest admissible
reconciliation governs the update, not the shortest drawn route. The grammar
of path cost has the same form. A direct displacement may look short, but
the cost belongs to the reconciliation the route must actually spend.

The Sagnac example supplies the cyclic face. Two routes can traverse the
same boundary in opposite directions and return to comparable endpoints while
accumulating unequal refinement tallies. The difference is not local
distance. It is the holonomy of the bookkeeping rule. Here again the grammar
separates endpoint relation from path cost. A closed cycle can return to the
same reference while still charging a different cyclic debt.

For the alpha reading, this separation has a precise consequence. The
dimensionless result may be recovered from the signed charge relation after
the three-rung aperture and Richardson bound have done their work. But the
recovery is not complete until the path cost has been read as well. A direct
and a composed route that agree in signed displacement may support the same
\(\alpha\) estimate while certifying different costs of obtaining it.

The grammar therefore gives two verdicts, not one. On the displacement side,
it permits a quotient: the composed and direct readings identify when their
signed residue depth agrees. On the cost side, it carries the difference:
the route through the intermediate rung and the route that bypasses it are
not the same expenditure unless a separate cost bound forces them to be
indistinguishable. The quotient of displacement never erases the obligation
to read cost.

This also prevents a false contradiction. If two routes produce the same
signed displacement at different cost, the grammar has not failed. It has
learned more. It has read that the endpoint relation is invariant under the
composition, while the path expenditure is not. The invariant belongs to the
alpha estimate. The expenditure belongs to the finite process by which the
estimate was made.

Thus the path relation completes the chapter's core. Seed and residue make
the repeated trial one history. The three rungs make the aperture readable.
Depth doubling and sign-flipped error make the extrapolation bounded. The
path relation identifies signed displacement while preserving cost. Alpha
enters only after all four obligations have been carried.

\section{What Curvature Carries Forward}

The chapter has not installed \(\alpha\) as a primitive constant. It has
forced a finite reading. A repeated trial carried seed and residue. A bounded
aperture supplied three rungs: charge, mass, and value. Depth doubling read a
sign-flipped sampling error and bounded the remaining residue. The path
relation identified signed displacement while preserving route cost. Only
after those obligations were in place could \(\alpha\) be carried forward as
an extrapolated dimensionless receipt.

This receipt is finite in ancestry. It remembers the trial that produced it:
the seed by which the next sample was selected, the residue restored after
each pass, the rung through which the aperture was read, the depth pair that
licensed extrapolation, and the path cost that was not erased by endpoint
agreement. The receipt may be used as a value in the next comparison, but it
does not cease to be answerable to that history.

The chapter has also kept charge and mass apart without separating them into
different worlds. Charge supplied the signed loop count. Mass supplied the
strain scale. Value supplied the comparable finite result. Alpha read their
relation inside one aperture. The grammar therefore avoids both false
simplicity and false proliferation. It does not reduce alpha to charge alone,
and it does not invent a fourth independent rung.

The Richardson-style move adds another discipline. The extrapolated reading
is sharper than either finite depth, but only because the paired depths have
made the leading error visible as a signed quantity. If the sign-flip is not
read, extrapolation is not forced. If the remainder is not bounded, the
improvement is only a hope. The alpha receipt carries the bound with the
value.

The path relation supplies the final accounting. The same signed displacement
can arrive by a direct route or a composed route. The grammar may quotient
the displacement when the residues agree, but it must still carry the cost
of the route. A finite alpha receipt therefore records not merely what value
was read, but how the reading was paid for.

Chapter 10 receives exactly this kind of object. A horizon or opaque surface
will not permit every interior distinction to pass through. It will export
only the readings the boundary can carry. Because this chapter has made
\(\alpha\) a bounded finite receipt, the next chapter can ask how it travels
with mass, charge, and spin-gravity residue at the surface. Alpha is no
longer a hidden interior possession. It is a carried surface-ready reading,
with its seed, residue, rung, extrapolation bound, and path cost still
attached.
